\section{Yesterday as evidence: rows and the gate}
\label{sec:infgate}

We know where the data is silent.  This section builds the mechanism
that lets yesterday fill the silence --- and proves that the
mechanism cannot touch anything the data has measured.

\subsection{Two estimates of one number}
\label{sec:infcombine}

Start with the smallest possible case, because the whole chapter is
this case repeated.  Suppose two independent, unbiased estimates of
the same unknown quantity are on the table --- say yesterday's
(transported) reading of some smile feature, and today's data-only
reading --- with known variances $\mathcal{V}_1$ and $\mathcal{V}_2$.
(``Unbiased'' means each is right on average; ``variance'' is the
familiar mean squared error of an estimate around the truth.)  How
should they be combined?  Try a weighted average with weight $a$ on
the first,
\[
\text{combined}=a\cdot\text{first}+(1-a)\cdot\text{second}.
\]
Because the two are independent and unbiased, the combination is
unbiased for every $a$, and its variance is
\[
a^2\,\mathcal{V}_1+(1-a)^2\,\mathcal{V}_2 .
\]
Minimize over $a$: differentiate and set to zero,
\[
2a\,\mathcal{V}_1-2(1-a)\,\mathcal{V}_2=0
\qquad\Longrightarrow\qquad
a=\frac{\mathcal{V}_2}{\mathcal{V}_1+\mathcal{V}_2}.
\]
Each estimate is weighted by the \emph{other's} variance ---
equivalently, by its own \emph{precision}, the reciprocal
$1/\mathcal{V}$ of its variance: dividing numerator and denominator
by $\mathcal{V}_1\mathcal{V}_2$,
\[
a=\frac{\mathcal{V}_2}{\mathcal{V}_1+\mathcal{V}_2}
=\frac{1/\mathcal{V}_1}{1/\mathcal{V}_1+1/\mathcal{V}_2}.
\]  Substituting the optimal $a$, the combined variance is
\begin{equation}
\mathcal{V}_{\mathrm{comb}}
=\frac{\mathcal{V}_1\,\mathcal{V}_2}{\mathcal{V}_1+\mathcal{V}_2},
\qquad\text{i.e.}\qquad
\frac{1}{\mathcal{V}_{\mathrm{comb}}}
=\frac{1}{\mathcal{V}_1}+\frac{1}{\mathcal{V}_2}:
\label{eq:infprecadd}
\end{equation}
precisions add.  A worked number to keep: yesterday says $20.0\%$
with standard deviation $0.30\%$; today's fit says $20.4\%$ with
standard deviation $0.15\%$.  Variances $0.09$ and $0.0225$ (in
$\%^2$); weight on yesterday $0.0225/0.1125=0.20$; combination
$0.20\times20.0+0.80\times20.4=20.32\%$.  Today, four times as
precise in variance, gets $80\%$ of the say.  These particular
numbers return in \cref{sec:inffilter}.

This is ordinary Bayesian updating in its simplest clothing: the
first estimate plays the \emph{prior} (what was believed before
today's data), the second the data, and the combination is the
\emph{posterior} (the belief after).  The share of the posterior
owned by the prior is
\begin{equation}
\text{prior share}
=\frac{1/\mathcal{V}_1}{1/\mathcal{V}_1+1/\mathcal{V}_2},
\label{eq:infshare}
\end{equation}
its precision divided by the total.  Keep \cref{eq:infshare} in
sight; the whole design question of this section is what that share
should be when the ``data'' side is a whole fit rather than one
number.

\subsection{The prior as extra rows}

How does yesterday enter a real calibration?  Chapters 2--4 fitted a
parameter vector $\theta$ (the model's coordinates, in whatever
chart) by least squares over quote equations.  The prior joins the
same objective as \emph{pseudo-observations}: pick a set of smile
functionals $O_j$ --- a functional here is just a number read off a
smile, such as the fitted volatility at one strike, or a difference
of two such volatilities --- evaluate each on yesterday's transported
surface to get a target $\hat O_j^{-}$, and append one row per
functional:
\begin{equation}
\text{row } j:\qquad
\sqrt{\lambda_j}\,\bigl(O_j(\theta)-\hat O_j^{-}\bigr),
\label{eq:infrow}
\end{equation}
carried at weight $\lambda_j\ge0$ alongside the quote rows.  (The
minus superscript marks a value believed \emph{before} today's data;
it will do heavy duty from \cref{sec:inffilter} on.)  Squared and
summed with the data misfit, such a row is exactly the prior side of
\cref{eq:infshare} for the coordinate $O_j$, with prior precision
$\lambda_j$.  Everything therefore reduces to one question: what
should $\lambda_j$ be?

\subsection{The design fork: a ridge, or a gate}

The familiar answer is a \emph{ridge} --- Tikhonov regularization, a
fixed small $\lambda_j$ for every row, the standard cure for
ill-posed fitting (a problem is ill-posed when its data do not
determine a unique, stable answer; Chapter~4's extraction was one,
and Chapter~4 used exactly this cure on the local-volatility sheet)
\citep{Tikhonov1977}.  For the present purpose the ridge has a
disqualifying property.  By \cref{eq:infshare}, a fixed
$\lambda>0$ owns the share $\lambda/(\lambda+\mathcal{I})$ of
\emph{every} coordinate, where $\mathcal{I}$ is the precision with
which today's data pins that coordinate.  The share shrinks as the
data improves --- but it never reaches zero.  A ridge biases the
superbly-quoted at-the-money level toward yesterday exactly as
surely as it stabilizes the dark wing, only less.  If yesterday's
level was $20\%$ and today's true level is $24\%$, a ridge pulls
today's published level below $24\%$ --- it \emph{damps} a genuine
move the market has plainly made.  For a fitter whose first duty is
to report what the market says, that is not a small flaw; it rules
the ridge out.

What is wanted instead is a weight that depends on the data: full
strength where the data is silent, and \emph{exactly zero} --- not
merely small --- once the data identifies the coordinate.  The
reference implementation uses one function for this everywhere:

\begin{definition}[Activation gate]
\label{def:infgate}
For an identification precision $\mathcal{I}\ge0$, a required
precision $\mathcal{I}_{\mathrm{req}}>0$, and a sharpness exponent
$\gamma>0$,
\begin{equation}
\operatorname{gate}(\mathcal{I})
=\Bigl(\min\bigl(\max\bigl(1-\mathcal{I}/\mathcal{I}_{\mathrm{req}},
\,0\bigr),\,1\bigr)\Bigr)^{\gamma}.
\label{eq:infgate}
\end{equation}
\end{definition}

Read it from the inside out: $1-\mathcal{I}/\mathcal{I}_{\mathrm{req}}$
is the shortfall of precision below the requirement; the
$\max(\cdot,0)$ clips it at zero once the requirement is met; the
$\min(\cdot,1)$ caps it at one when there is no data at all; and the
exponent $\gamma$ shapes how fast the gate closes as data
accumulates.  Each persisted row then receives
$\lambda_j\propto\operatorname{gate}(\mathcal{I}_j)$: a fixed
budget of prior weight is split across the rows in proportion to
their gate values, and a row whose gate is zero is simply dropped.

\begin{proposition}[The no-damp guarantee]
\label{prop:infnodamp}
If coordinate $j$ satisfies
$\mathcal{I}_j\ge\mathcal{I}_{\mathrm{req}}$, then
$\operatorname{gate}(\mathcal{I}_j)=0$, hence $\lambda_j=0$, and row
$j$ of \cref{eq:infrow} is the zero function of $\theta$.  The
objective, its gradient, and every downstream quantity are then
independent of the prior value $\hat O_j^{-}$.  In particular, when
every gate is closed the fit is the prior-free fit.
\end{proposition}

\begin{proof}
With $\lambda_j=0$ the row is $\sqrt{0}\cdot(O_j(\theta)-\hat
O_j^{-})=0$ identically: it contributes zero to the objective and a
zero row to the derivative stack, for every $\theta$.  Nothing that
is computed from the objective can depend on $\hat O_j^{-}$ through
a row that is identically zero.
\end{proof}

The proof is three lines because the content is structural, not
analytic: the gate does not make the prior's influence small where
data speaks --- it makes the prior \emph{absent} there.  That is the
difference between gating and shrinkage, and it turns ``do not damp
the signal'' from a preference into a property of the objective.

\begin{figure}[htbp]
\centering
\paperfig{\textwidth}{fig_flt_gate.pdf}
\caption{The gate, and why a ridge cannot replace it.  (a)~The gate
\cref{eq:infgate} for three sharpness exponents, against precision in
units of the requirement: all three reach exactly zero at the dotted
requirement line and stay there --- the dead zone.  (b)~The share of
the posterior owned by the prior, \cref{eq:infshare}, with prior
weight $\lambda$ fixed for the ridge (dashed) and gated (solid), both
at unit budget: at twice the required precision the ridge still owns
\MacPriorGateRidgeSharePct\% of the coordinate; the gated prior owns
exactly nothing.}
\label{fig:infgate}
\end{figure}

\Cref{fig:infgate} draws both halves of the argument.  In
panel~(a), the three curves are the gate at $\gamma=0.5$, $1$, and
$2$; all three start at one where the data says nothing, descend as
precision builds, and hit zero exactly at the required precision ---
after which they stay at zero, which is the whole point.  The
exponent only shapes the descent: $\gamma=2$ releases the prior
early, $\gamma=0.5$ holds on longer.  Panel~(b) translates gates
into posterior shares via \cref{eq:infshare}.  The dashed curve is
the ridge: at a precision of twice the requirement --- a coordinate
the data pins twice as well as demanded --- the ridge still owns
\MacPriorGateRidgeSharePct\% of the answer.  The solid curve is the
gated prior: identical in the dark, and exactly zero past the
requirement.  Same stabilization where it is wanted, none where it
is not.

We now have the mechanism and its guarantee.  What the gate consumes
is a precision $\mathcal{I}_j$ per persisted functional --- a number
that must be produced cheaply, before fitting, for functionals more
interesting than a single strike.  Producing it is the next
section's job.
